% Printed 498; principal PDF 505; Google PDF 509.
\par\clearpage
\noindent\makebox[0pt][l]{\fontsize{13}{16}\selectfont 498}\hfill{\small\textbf{CORRESPONDENCE.}}\hfill\null\par\medskip
\noindent that at my observatory no really suitable setup can be obtained for such trial measurements, that few suitable objects are visible, and finally that winter is unfavourable to my very sensitive constitution. I will at least try what is within my powers, however.

\par\noindent\hspace*{\parindent}If you should happen to have had similar experiences yourself and undertaken various related trials as a result, you would much oblige me by communicating them. It is indeed a vital question for theodolites of this construction. In \textsc{Borda} circles the error cannot occur at all by virtue of the principle of construction. There is no reason why moving the alidade should displace the lower telescope \textls[120]{relative to the graduated} circle, and moving the whole circle likewise does not displace the alidade relative to the graduated circle. Theodolites can be built on the same principle, but then they are quite different instruments from \textsc{Reichenbach}'s. Unless I am mistaken, \textsc{Schumacher} has one of this kind by \textsc{Gambey} (in which case the lower telescope plays a quite different role), but I do not know the details of the arrangement.\enspace .\ .\ .\ .\ .\ .

\par\noindent\hspace*{\parindent}It is striking how things put down on paper make a different impression from the one they make when one only thinks them through. On looking over the preceding lines again quickly, it seems to me that one circumstance must appear more important to an impartial reader who has had no experience of his own than I myself have hitherto thought it. This is that \textls[120]{at first} none of the theodolites showed the phenomenon at all, or at least not perceptibly. Some kind of \textls[120]{wear} must therefore be at work, and I will probably have new screws, nuts, and balls made in any case, even though, as I said, no play can be detected in these parts individually. I will try, however, whether a \textls[120]{much more powerful} telescope might already reveal something of such play in individual parts.

\par\bigskip
\begin{center}{\large\textsc{Gauss} to \textsc{Bessel}.\quad Göttingen, 15 August 1844.}\end{center}
\par\noindent\hspace*{\parindent}.\ .\ .\ .\ .\ .\enspace The subject on which I last corresponded with you, namely the tendency of repeating \textls[120]{theodolites} to give angles

% Printed 499; principal PDF 506; Google PDF 510.
\par\clearpage
\noindent\null\hfill{\small\textbf{MEASUREMENT ERRORS.}}\hfill\makebox[0pt][r]{\fontsize{13}{16}\selectfont 499}\par\medskip
\noindent that are too small, continued to trouble me greatly last winter. The cause can only be that the circle does not remain absolutely fixed during those manipulations of the alidade in which it is assumed to remain fixed. The alterations I had made have, however, proved very effective, both in the extremely numerous trial measurements I made myself from January through March and in the actual use my son continues to make of the instrument. Some of last year's results were so poor that they had to be discarded altogether, whereas this year's are all as good as can be expected from an instrument of this size. The main alteration is that the alidade pivot and its bushing (which in turn forms the pivot of the graduated circle) are each supported by a spring so arranged that either supporting spring can be tensioned or slackened independently of the other at will and at once (wholly or in part). In use there is then a constant \textls[120]{systematic} alternation between the states of these springs, so that, counting everything, 19 operations are required for each repetition (including the reading). Yet one learns the sequence of these operations so mechanically that one can work almost as quickly as with the usual arrangement.
\par\bigskip
\begin{center}\rule{2.6cm}{.4pt}\end{center}
\vspace{1.2cm}
{\small
\begin{center}REMARKS.\end{center}
\par\noindent\hspace*{\parindent}The note on the instruments used in the Hanover arc measurement and land survey, and on those still required for the latter, is taken from a report by \textsc{Gauss} of 21 March 1829 to Privy Cabinet Councillor \textsc{Hoppenstedt}, held in the Gauss archive.
\par\noindent\hspace*{\parindent}In connection with the preceding letters to \textsc{Olbers}, \textsc{Schumacher}, and \textsc{Bessel}, printed from the originals, compare the station adjustment for Brillit, pp. 265/267, and a remark by Professor \textsc{Goldschmidt}, p. 289.
\par\noindent\hfill\textsc{Krüger.}\hspace*{1em}\par}
\vspace{1cm}
\begin{center}\rule{3.5cm}{.4pt}\end{center}
\par\medskip\noindent\makebox[.86\linewidth][r]{63\textsuperscript{*}}\par

% Printed 500; principal PDF 507; Google PDF 511.
\par\clearpage
\noindent\makebox[0pt][l]{\fontsize{13}{16}\selectfont 500}\hfill{\small\textbf{CORRESPONDENCE.}}\hfill\null\par
\vspace*{3.2cm}
\begin{center}
{\LARGE CORRESPONDENCE AND POSTHUMOUS PAPERS.}\par
\vspace{1.5cm}
{\large [Calculation of the mean reading error including graduation error\par in the theodolite.]}\par
\bigskip\bigskip
{\large [1.]}\par\medskip
{\large\textsc{Gauss} to \textsc{Gerling}.\quad Göttingen, 17 April 1844.}
\end{center}
\par\noindent\hspace*{\parindent}.\ .\ .\ .\ .\ .\enspace I must reserve for another occasion writing to you about another check on theodolite observations by calculation (the purpose of which is to determine the mean error of the \textls[120]{readings}, including the graduation error). I made these calculations 19 years ago already, and they gave me much pleasure. For \textls[120]{every} complete set of readings that occurred I calculated $a-b+c-d=l$, where (omitting the degrees) $a$, $b$, $c$, $d$ denote the 4 readings. This $l$ ought (\textls[120]{despite} the eccentricity, indeed despite any variation in eccentricity) to be \textls[120]{constant}. It varies only as a result of errors of reading and graduation. If we put the mean value from very many, $n$, readings $=\lambda$ (which can be found with great precision for a given theodolite from many observations and ought to be invariable), and put $\frac{\Sigma(l-\lambda)^2}{n-1}=kk$, then $\frac12k$ is the mean error to which one is liable in reading one index, and $\frac14k$ the mean error to which one is liable in taking the mean of all 4 indices. Now I find the latter, or $\frac14k$, from complete sets of readings:

% Printed 501; principal PDF 508; Google PDF 512.
\par\clearpage
\noindent\null\hfill{\small\textbf{MEASUREMENT ERRORS.}}\hfill\makebox[0pt][r]{\fontsize{13}{16}\selectfont 501}\par\medskip
\begin{center}
\begingroup\renewcommand{\arraystretch}{1.3}\setlength{\tabcolsep}{11pt}
\begin{tabular}{@{\vrule width .8pt\hspace{11pt}}c|c|c@{\hspace{11pt}\vrule width .8pt}}
\noalign{\hrule height .8pt}
[Year & No. of readings & Mean error]\\[4pt]
\hline
1821 & \makebox[2.4em][r]{410} & \makebox[3em][r]{$1\overset{\prime\prime}{,}25$}\\[7pt]
1822 & \makebox[2.4em][r]{717} & \makebox[3em][r]{$1{,}57$}\\
1823 & \makebox[2.4em][r]{449} & \makebox[3em][r]{$1{,}61$}\\
1824 & \makebox[2.4em][r]{843} & \makebox[3em][r]{$1{,}60$}\\
1825 & \makebox[2.4em][r]{609} & \makebox[3em][r]{$1{,}55$}\\
\hline
1822---1825 & \makebox[2.4em][r]{2618} & \makebox[3em][r]{$1{,}58$}\\
\end{tabular}
\endgroup
\end{center}
\par\noindent\hspace*{\parindent}The very close agreement among the years 1822---1825 already suggests that the divergent value for 1821 is not accidental. Indeed, the graduation lines on \textsc{Schumacher}'s theodolite were a little finer than on my own. His dated from the time when \textsc{Reichenbach} still had the workshop himself. Mine was ordered from \textsc{Reichenbach}, but was made by \textsc{Ertel} under his supervision. On my 8-inch instrument (of 1829) the graduation is also appreciably less fine than on the 8-inch instrument by \textsc{Reichenbach} that I acquired in 1813. The former, however, has a \textls[120]{much} better .telescope.

\par\noindent\hspace*{\parindent}I see that after all I have now written to you about the most essential points of this subject, and so merely note, for completeness, a substantial shortcut to that calculation. For $\Sigma l=n\lambda$, hence
\[
\Sigma(l-\lambda)^2=\Sigma ll-2n\lambda\lambda+n\lambda\lambda=\Sigma ll-n\lambda\lambda
\left[=\Sigma ll-\frac{(\Sigma l)^2}{n}\right],
\]
\noindent hence
\[
k=\frac{\Sigma ll}{n-1}-\frac{n}{n-1}\lambda\lambda
\left[=\frac{\Sigma ll}{n-1}-\frac{(\Sigma l)^2}{n(n-1)}\right].
\]
\par\bigskip\bigskip
\begin{center}{\large [2.]}\end{center}
\par\noindent\hspace*{\parindent}To ascertain the errors arising from graduation and reading, all complete sets of readings on the theodolite, 449 in number [from the year 1823], were analysed as follows.

% Printed 502; principal PDF 509; Google PDF 513.
\par\clearpage
\noindent\makebox[0pt][l]{\fontsize{13}{16}\selectfont 502}\hfill{\small\textbf{POSTHUMOUS PAPERS.}}\hfill\null\par\medskip
\par\noindent\hspace*{\parindent}Sum of the readings at the first and third indices minus the sum of the readings at the second and fourth $=l$.
\[
\begin{aligned}
\Sigma l&=+133\\
\Sigma ll&=\phantom{+}18521,
\end{aligned}
\]
\noindent hence the mean value of $l=\left[\frac{133}{449}=\right]+0\overset{\prime\prime}{,}30=\lambda$;
\[
\text{mean value of }(l-\lambda)^2=
\left[\frac{1}{448}\left(18521-\frac{133^2}{449}\right)=\right]41{,}254.
\]
\par\noindent\hspace*{\parindent}Thus the mean error with one vernier, including graduation error:
\[
\left[\frac{\sqrt{41{,}254}}{2}=\right]\pm3\overset{\prime\prime}{,}211.
\]
\par\noindent\hspace*{\parindent}On \textsc{Schumacher}'s theodolite, 410 readings from 1821 gave
\[
\begin{aligned}
\Sigma l&=-1896\\
\Sigma ll&=\phantom{-}18982.
\end{aligned}
\]
\par\noindent\hspace*{\parindent}Mean value of $l=-4\overset{\prime\prime}{,}62=\lambda$.
\par\noindent\hspace*{\parindent}Mean value of $(l-\lambda)^2=24{,}97$.
\par\noindent\hspace*{\parindent}Mean reading error of one vernier including graduation error $=\pm2\overset{\prime\prime}{,}499$.
\par\bigskip\bigskip
\par\noindent\hspace*{\parindent}In 1822 the theodolite gave, from 717 readings,
\[
\begin{aligned}
\Sigma l&=+212\\
\Sigma ll&=\phantom{+}28433.
\end{aligned}
\]
\par\noindent\hspace*{\parindent}Mean value of $l=+0\overset{\prime\prime}{,}30=\lambda$.
\par\noindent\hspace*{\parindent}Mean value of $(l-\lambda)^2=39{,}623$.
\par\noindent\hspace*{\parindent}Mean reading error with one vernier including graduation error $=\pm3\overset{\prime\prime}{,}147$.
\par\bigskip\bigskip
\par\noindent\hspace*{\parindent}This result agrees so closely with that of 1823 that we can combine the two. Thus from 1166 readings we have:

% Printed 503; principal PDF 510; Google PDF 514.
\par\clearpage
\noindent\null\hfill{\small\textbf{MEASUREMENT ERRORS.}}\hfill\makebox[0pt][r]{\fontsize{13}{16}\selectfont 503}\par\medskip
\[
\begin{aligned}
\Sigma l&=+345\\
\Sigma ll&=\phantom{+}46954\\
\lambda&=+0\overset{\prime\prime}{,}296.
\end{aligned}
\]
\par\noindent\hspace*{\parindent}Mean value of $(l-\lambda)^2=40{,}216$.
\par\noindent\hspace*{\parindent}Mean error with one vernier $=\pm3\overset{\prime\prime}{,}1708$.
\par\noindent\hspace*{\parindent}The frequencies of the individual values of $l$, disregarding the sign, were as follows:
\begin{center}
\begingroup\renewcommand{\arraystretch}{1.15}\setlength{\tabcolsep}{10pt}
\begin{tabular}{@{\vrule width .8pt\hspace{10pt}}c|c@{\hspace{10pt}\vrule width .8pt\hspace{10pt}}c|c@{\hspace{10pt}\vrule width .8pt\hspace{10pt}}c|c@{\hspace{10pt}\vrule width .8pt}}
\noalign{\hrule height .8pt}
\makebox[0pt][l]{$[$}\makebox[2em]{$l$} & Count & \makebox[2em]{$l$} & Count & \makebox[2em]{$l$} & Count $]$\\[4pt]
\hline
\makebox[.6em][r]{0}\rlap{$''$} & \makebox[1.8em][r]{45} & \makebox[1.2em][r]{6}\rlap{$''$} & \makebox[1.8em][r]{133} & \makebox[1.2em][r]{11}\rlap{$''$} & \makebox[1.2em][r]{38}\\
\makebox[.6em][r]{1} & \makebox[1.8em][r]{106} & \makebox[1.2em][r]{7} & \makebox[1.8em][r]{104} & \makebox[1.2em][r]{12} & \makebox[1.2em][r]{24}\\
\makebox[.6em][r]{2} & \makebox[1.8em][r]{108} & \makebox[1.2em][r]{8} & \makebox[1.8em][r]{82} & \makebox[1.2em][r]{13} & \makebox[1.2em][r]{15}\\
\makebox[.6em][r]{3} & \makebox[1.8em][r]{123} & \makebox[1.2em][r]{9} & \makebox[1.8em][r]{84} & \makebox[1.2em][r]{14} & \makebox[1.2em][r]{9}\\
\makebox[.6em][r]{4} & \makebox[1.8em][r]{113} & \makebox[1.2em][r]{10} & \makebox[1.8em][r]{52} & \makebox[1.2em][r]{15} & \makebox[1.2em][r]{4}\\
\makebox[.6em][r]{5} & \makebox[1.8em][r]{126} &  &  &  & \\
\end{tabular}
\endgroup
\end{center}
\par\noindent\hspace*{\parindent}In 1824 the theodolite gave, from 843 readings,
\[
\begin{aligned}
\Sigma l&=+104\\
\Sigma ll&=\phantom{+}34506\\
\lambda&=+0\overset{\prime\prime}{,}12.
\end{aligned}
\]
\par\noindent\hspace*{\parindent}Mean value of $(l-\lambda)^2=40{,}966$.
\par\noindent\hspace*{\parindent}Mean error with one vernier including graduation error $=\pm3\overset{\prime\prime}{,}200$.
\par\noindent\hspace*{\parindent}Thus from all 3 years [1822---1824], 2009 readings,
\[
\begin{aligned}
\Sigma l&=+449\\
\Sigma ll&=\phantom{+}81460\\
\lambda&=+0\overset{\prime\prime}{,}2235.
\end{aligned}
\]
\par\noindent\hspace*{\parindent}Mean value of $(l-\lambda)^2=40{,}518$.
\par\noindent\hspace*{\parindent}Mean error of one vernier $=\pm3\overset{\prime\prime}{,}1827$. The frequencies of the individual values of $l$, disregarding the sign, were as follows:

% Printed 504; principal PDF 511; Google PDF 515.
\par\clearpage
\noindent\makebox[0pt][l]{\fontsize{13}{16}\selectfont 504}\hfill{\small\textbf{POSTHUMOUS PAPERS.}}\hfill\null\par\medskip
\begin{center}
\begingroup\renewcommand{\arraystretch}{1.15}\setlength{\tabcolsep}{10pt}
\begin{tabular}{@{\vrule width .8pt\hspace{10pt}}c|c@{\hspace{10pt}\vrule width .8pt\hspace{10pt}}c|c@{\hspace{10pt}\vrule width .8pt\hspace{10pt}}c|c@{\hspace{10pt}\vrule width .8pt}}
\noalign{\hrule height .8pt}
\makebox[0pt][l]{$[$}\makebox[2em]{$l$} & Count & \makebox[2em]{$l$} & Count & \makebox[2em]{$l$} & Count $]$\\[4pt]
\hline
\makebox[.6em][r]{0}\rlap{$''$} & \makebox[1.8em][r]{80} & \makebox[1.2em][r]{7}\rlap{$''$} & \makebox[1.8em][r]{191} & \makebox[1.2em][r]{14}\rlap{$''$} & \makebox[1.2em][r]{13}\\
\makebox[.6em][r]{1} & \makebox[1.8em][r]{182} & \makebox[1.2em][r]{8} & \makebox[1.8em][r]{155} & \makebox[1.2em][r]{15} & \makebox[1.2em][r]{6}\\
\makebox[.6em][r]{2} & \makebox[1.8em][r]{185} & \makebox[1.2em][r]{9} & \makebox[1.8em][r]{138} & \makebox[1.2em][r]{16} & \makebox[1.2em][r]{1}\\
\makebox[.6em][r]{3} & \makebox[1.8em][r]{195} & \makebox[1.2em][r]{10} & \makebox[1.8em][r]{101} & \makebox[1.2em][r]{17} & \makebox[1.2em][r]{1}\\
\makebox[.6em][r]{4} & \makebox[1.8em][r]{205} & \makebox[1.2em][r]{11} & \makebox[1.8em][r]{63} & \makebox[1.2em][r]{18} & \makebox[1.2em][r]{---}\\
\makebox[.6em][r]{5} & \makebox[1.8em][r]{217} & \makebox[1.2em][r]{12} & \makebox[1.8em][r]{38} & \makebox[1.2em][r]{19} & \makebox[1.2em][r]{1}\\
\makebox[.6em][r]{6} & \makebox[1.8em][r]{214} & \makebox[1.2em][r]{13} & \makebox[1.8em][r]{23} &  & \\
\end{tabular}
\endgroup
\end{center}
\begin{center}\rule{3.5cm}{.4pt}\end{center}
\par\bigskip
\par\noindent\hspace*{\parindent}On the [\textsc{Borda}] circle, in [1823], 46 complete sets of readings were taken. The results were
\[
\begin{aligned}
\Sigma l&=-174\\
\Sigma ll&=\phantom{-}1534\\
\lambda&=-3\overset{\prime\prime}{,}78.
\end{aligned}
\]
\par\noindent\hspace*{\parindent}Mean value of $(l-\lambda)^2=19{,}46$.
\par\noindent\hspace*{\parindent}Mean reading error of one vernier including graduation error $=\pm2\overset{\prime\prime}{,}206$.
\par\noindent\hspace*{\parindent}In 1821 the circle gave, with 92 readings,
\[
\begin{aligned}
\Sigma l&=-650\\
\Sigma ll&=\phantom{-}11578\\
\lambda&=-7\overset{\prime\prime}{,}07.
\end{aligned}
\]
\par\noindent\hspace*{\parindent}Mean value of $(l-\lambda)^2=76{,}767$.
\par\noindent\hspace*{\parindent}Mean reading error of one vernier including graduation error $=\pm4\overset{\prime\prime}{,}381$.
\par\noindent\hspace*{\parindent}In 1822, 137 readings were taken on the circle. These gave
\[
\begin{aligned}
\Sigma l&=-543\\
\Sigma ll&=\phantom{-}6155\\
\lambda&=-3\overset{\prime\prime}{,}96.
\end{aligned}
\]
\par\noindent\hspace*{\parindent}Mean value of $(l-\lambda)^2=29{,}432$.

% Printed 505; principal PDF 512; Google PDF 516.
\par\clearpage
\noindent\null\hfill{\small\textbf{MEASUREMENT ERRORS.}}\hfill\makebox[0pt][r]{\fontsize{13}{16}\selectfont 505}\par\medskip
\par\noindent\hspace*{\parindent}Mean reading error with one vernier including graduation error $=\pm2\overset{\prime\prime}{,}713$.
\par\noindent\hspace*{\parindent}If one allows oneself to combine the measurements of 1822 and 1823, one obtains [with 183 readings]
\[
\begin{aligned}
\Sigma l&=-717\\
\Sigma ll&=\phantom{-}7689\\
\lambda&=-3\overset{\prime\prime}{,}92.
\end{aligned}
\]
\par\noindent\hspace*{\parindent}Mean value of $(l-\lambda)^2=\left[26{,}812.\right.$
\par\noindent\hspace*{\parindent}Mean reading error with one vernier including graduation error $=\pm2\overset{\prime\prime}{,}589.]$
\par\bigskip\bigskip
\begin{center}\rule{2.6cm}{.4pt}\end{center}
\vspace{1.5cm}
{\small
\begin{center}REMARKS.\end{center}
\par\noindent\hspace*{\parindent}The preceding letter to \textsc{Gerling}, Art. [1], was printed from the original. Art. [2] was taken from an observation notebook for the Hanover arc measurement.
\par\noindent\hfill\textsc{Krüger.}\hspace*{1em}\par}
\vspace{2cm}
\begin{center}\rule{3.5cm}{.4pt}\end{center}

\vfill
\noindent IX.\hfill 64\hspace*{1cm}\par

% Contextual printed 506; principal PDF 513; Google PDF 517: unnumbered blank leaf.
\par\clearpage
\null

% Contextual printed 507; principal PDF 514; Google PDF 518: unnumbered section title.
\par\clearpage
\vspace*{4cm}
\begin{center}
{\fontsize{22}{30}\selectfont
REDUCTION OF OBLIQUE\par
ANGLES TO THE HORIZON.\par
OBSERVATIONS BY REPETITION.\par}
\end{center}
\vspace{8.2cm}
\noindent\makebox[.86\linewidth][r]{64\textsuperscript{*}}\par
